\chapter{Probabilistic Modelling}
\section{Probabilistic models}\label{sec:prob}

The risk analysis (Sections~3.1 and~3.2) needs the probability of failure of
each of the 17 walmuren that still stand on timber piles
(Section~\ref{sec:data}). We model it with two limit states, evaluated by
Monte Carlo simulation, and check the model against data and observations.
The full equations are in Appendix~\ref{app:prob}.

\subsection{Limit states}\label{sec:prob-ls}

A wall fails when its body (wall, floor and the soil on it) overturns about
the front pile row, or when a timber pile breaks in bending
\cite{hemel2023,hemel2024}:
\begin{align}
g_1 &= \theta_R M_R + M_\delta - S, \label{eq:g1}\\
g_2 &= \text{MOR} - \left(\theta_M \frac{32 M}{\pi d^3} + \frac{4N}{\pi d^2}\right),
\qquad d = D_0 - 2\bigl(s_0 + r\,(t - 2026)\bigr). \label{eq:g2}
\end{align}
$S$ is the moment of soil, water and surcharge, $M_R$ that of the weight of
the body and $M_\delta$ that of wall friction. $M$ and $N$ act in a pile whose
sound core $d$ shrinks as the rotted shell $s_0$ grows at rate $r$. The wall
fails if either limit state is violated (series system):
$P_f = P(g_1 < 0 \cup g_2 < 0)$ and $\beta = -\Phi^{-1}(P_f)$, estimated as
the fraction of failed draws (400{,}000 per wall and type; 2026 and 2041).

\subsection{Random variables and distributions}\label{sec:prob-variables}

Table~\ref{tab:prob-variables} lists the random variables. Each type of
distribution follows from what is known: a \emph{lognormal} for positive
quantities with a multiplicative error, such as $K$ and the model factors
$\theta$ \cite{jcss2001}; a \emph{normal} for unit weights and timber
properties, where a CoV of at most 0.20 makes negative values negligible; a
\emph{uniform} where only a range is known, as it adds nothing beyond the
bounds; and a \emph{Bernoulli} for the wall type.

\begin{table}[htbp]
\centering
\small
\caption{Random variables of the model. Per wall: values differ between
walls. CoV = coefficient of variation.}
\label{tab:prob-variables}
\begin{tabularx}{\textwidth}{@{}l l l >{\raggedright\arraybackslash}X@{}}
\toprule
Variable & Distribution & Parameters & Estimation \\
\midrule
Werf level (per wall) & Normal & AHN median; $\sigma = 0.05$\,m & AHN data \\
Wall type $T$ (per wall) & Bernoulli & $P(T{=}2)$ = 0.2--0.9 per rak & Dive notes \cite{tno2009} \\
Earth pressure $K$ & Lognormal & 0.35--0.60; CoV 0.20--0.25 & Moments per soil class \\
Soil unit weight $\gamma$ & Normal & 15--18\,kN/m$^3$; CoV 0.05--0.06 & Moments per soil class \\
Surcharge $q$; distance $a$ & Normal; uniform & 7\,kPa, CoV 0.20; 0--3\,m & \cite{hemel2024}; range \\
Groundwater & Uniform & +0.5 to +0.9\,m NAP & Range around canal level \\
Unit weights wall & Normal & 19 / 23\,kN/m$^3$; CoV 0.05 & Typical values \\
Floor level $\Delta z$ & Uniform & $\pm 0.15$\,m & Drawing accuracy \\
Model factors $\theta_R$, $\theta_M$ & Lognormal & mean 1; CoV 0.15, 0.10 & Judgement; JCSS \cite{jcss2001} \\
Pile diameter $D_0$ & Normal & 180\,mm; CoV 0.10 & \cite{ibu2010,hemel2023} \\
Pile spacing $s$ & Uniform & 0.8--1.2\,m & \cite{ibu2010} \\
Soft shell $s_0$; growth $r$ & Uniform & 0--45\,mm; 0.10--0.25\,mm/yr & Measured range \cite{hemel2023} \\
MOR; $E$ & Normal & 22\,MPa, CoV 0.10; 9\,GPa, 0.20 & \cite{hemel2023} \\
Subgrade modulus $k$ & Uniform & 4{,}300--10{,}000\,kN/m$^3$ & \cite{hemel2023} \\
\bottomrule
\end{tabularx}
\end{table}

The parameters are estimated in three ways: from a literature mean and CoV
(method of moments), from our own data (werf level, wall-type prior), and by
Bayesian updating with the observation that every wall still stands (proven
strength \cite{hemel2024}). Draws in which the wall would already fail today
under soil and water alone are rejected; the same evidence updates
$P(T{=}2)$ (Appendix~\ref{app:prob}).

\subsection{Dependence between the variables}\label{sec:prob-dependence}

The dependence is modelled explicitly in four ways:
\begin{enumerate}\setlength{\itemsep}{0pt}\setlength{\parskip}{0pt}
  \item \textbf{Common cause.} The wall type sets the geometry, the floor
    level, the rake of the piles and the head fixity together. The soil class
    sets the means of $K$ and $\gamma$. $P_f$ is a mixture:
    $P_f = \sum_T P(T \mid E)\,P_f(T \mid E)$.
  \item \textbf{Shared variables.} One draw of $K$ enters the soil pressure,
    the surcharge pressure and the wall friction ($\delta = \tfrac23\varphi$,
    with $\varphi$ from $K$). One groundwater level enters the soil weight,
    the water pressure and the buoyancy of the body.
  \item \textbf{Same draws.} Both limit states use the same draws, so their
    correlation is kept: for type 2, overturning and pile failure occur
    together 3.4 times as often as if they were independent. The years 2026
    and 2041 share the piles; only the shell grows.
  \item \textbf{Proven strength.} The check ``stands today'' shares all
    resistance and soil variables with the future load case, but has its own
    groundwater and no surcharge.
\end{enumerate}
Other variables are independent, as no data on their correlation exist. We
tested the two most likely correlations (Appendix~\ref{app:prob}): denser
soil is heavier with a lower $K$, $\rho(K,\gamma) = -0.5$, and stronger timber
is stiffer, $\rho(\text{MOR}, E) = 0.8$ \cite{jcss2001}. Each changes $\beta$
by at most 0.03, mostly upwards, so independence is acceptable and slightly
conservative. The walls are simulated independently, although in reality
they share the uncertain pile condition, surcharge and model factors; this
matters when walls are combined in Section~3.1.

\subsection{Model quality and validity}\label{sec:prob-validity}

Few inputs can be tested against data, but where they can:
\begin{itemize}\setlength{\itemsep}{0pt}\setlength{\parskip}{0pt}
  \item \textbf{Werf level.} AHN points on a werf scatter by about 0.08\,m,
    so the wall average is known to about 0.04\,m, close to $\sigma =
    0.05$\,m. For three walls (Rak 03 OZ, Kromme Nieuwe Gracht 002 and 009)
    the median includes points on stairs or the street, which makes their
    werf 0.3--0.5\,m too high and their $\beta$ too low.
  \item \textbf{Groundwater.} The BRO wells, interpolated to the walls, give
    +0.9 to +2.1\,m NAP, mostly above the werf itself: they stand under the
    raised streets and do not describe the werf, which drains to the canal.
    The range is therefore a physical bound, not a fit. Raising it to
    $U(+0.6, +1.0)$\,m NAP lowers $\beta$ by 0.04--0.09 (type 2) and 0.1--0.4
    (type 1).
  \item \textbf{Soil class.} Only 25--67\,\% of the borings within 150\,m
    give the class of the nearest boring. Treating the class as uncertain,
    with weights from these borings, changes $\beta$ by $-0.7$ to $+0.3$. Walls
    classed as sand lose most, as half of their nearby borings find clay or
    peat. This is the weakest part of the input model.
  \item \textbf{Pile shell.} $U(0, 45)$\,mm has a mean of 22.5\,mm, close to
    the 21\,mm measured in Amsterdam \cite{hemel2023}; no Utrecht data exist.
\end{itemize}

The strongest test is that all 17 walls stand, so a valid model must let
them stand in most draws. A first model that treated the wall as a masonry
block failed in almost every draw and was rejected. With the measured
sections, type 2 stands in 87--100\,\% of the draws and type 1 in only
10--98\,\%: the type-1 model is conservative, or many walls are type 2, and
the update shifts $P(T{=}2)$ accordingly. The Monte Carlo error is small
(CoV of $\hat P_f$ below 2\,\%, $\beta$ within $\pm 0.01$).

The model is valid for low walmuren with a werf on their old timber piles,
for the two Oudegracht cross-sections. It does not cover settlement, sliding
or failure of the timber screen. Because the main inputs are literature
values and ranges, not fits to Utrecht data, the results show the ranking and
order of magnitude of $P_f$ rather than exact values (Section~4.2).

\newpage
%% ----------------------------------------------------------------------
%%  MAX 2 PAGES
%%
%%  TASK DESCRIPTION
%%  To perform the risk analysis in Sections \ref{sec:fault-tree} and
%%  \ref{sec:trade-off}, you need probabilistic models. Identify the
%%  relevant variables and clarify the (in)dependence between them.



%%  Construct adequate probabilistic models with explicit consideration of
%%  the dependence between the variables. You may use uni- or
%%  multivariate probability distributions to describe or quantify the
%%  randomness in the data. You may apply different parameter estimation
%%  techniques.


%%  Justify the selection of the model(s) and evaluate model quality or
%%  validity. For example: does the model fit the data well? To what
%%  extent is the model valid to represent the data?
%%  Use the appendices to add supporting information.
%%




%%  CHECKLIST (challenge task 4: build probabilistic models of
%%  deterioration and demand)


%%   - Represent the relevant degradation processes and their drivers:
%%     groundwater fluctuation and drought acting on timber pile
%%     foundations, bacterial and fungal decay, corrosion, masonry
%%     degradation, scour, and pressure from tree roots.
%%   - Represent the future demands placed on the assets: urban
%%     densification, heavier vehicles and construction traffic, vessel
%%     impact and mooring, and a changing climate (KNMI'23 scenarios
%%     project drier summers, wetter winters and more extreme rainfall).
%%   - Treat dependence between assets and between drivers explicitly
%%     rather than assuming it away.
%%   - Consider that the assessment rules themselves are a source of
%%     uncertainty: they embed safety margins calibrated for new design,
%%     and the margin between calculated and actual capacity is itself an
%%     uncertain quantity worth estimating.
%% ----------------------------------------------------------------------


